% Lecteur de l'édition française avec compléments.
\documentclass[11pt,a4paper,openany]{memoir}
\usepackage{mathpazo}
\usepackage{fontspec}
\setmainfont{TeX Gyre Pagella}
\setsansfont{TeX Gyre Heros}
\setmonofont{TeX Gyre Cursor}
\usepackage{natbib}
\newcommand*{\olpath}{../..}
\newcommand*{\ollangid}{fr}
\newcommand*{\ollanguage}{french}
\input{\olpath/sty/open-logic.sty}
\input{\olpath/sty/open-logic-defer.sty}
\includeenv{editorial}
\usepackage{footnotehyper}
\makesavenoteenv{editorial}
\frcompletevariantstrue
\problemsperchapter
\setlrmarginsandblock{28mm}{28mm}{*}
\setulmarginsandblock{26mm}{28mm}{*}
\checkandfixthelayout
\linespread{1.05}
\raggedbottom
\setlength{\emergencystretch}{2em}
\hfuzz=3pt
\let\cleardoublepage\clearpage
\pagestyle{plain}
\setlength{\cftbeforechapterskip}{0.65em}
\hypersetup{pdftitle={OpenLogic — Édition française},
pdfauthor={Open Logic Project ; traduction française produite par IA},
pdfsubject={Logique, ensembles, calculabilité, théorie des modèles et méthodes formelles},
linkcolor=black,citecolor=black,urlcolor=blue}

\begin{document}
\frontmatter
\begin{titlingpage}
\setlength{\parindent}{0pt}
\vspace*{15mm}
{\sffamily\Huge\bfseries OpenLogic\par}
\vspace{6mm}
{\sffamily\LARGE Édition française\par}
\vspace{14mm}
\vfill
Ce volume réunit la théorie des ensembles, les logiques propositionnelle
et du premier ordre, la théorie des modèles, la calculabilité,
l'incomplétude, les logiques modales et non classiques, ainsi que les
méthodes de démonstration et des éléments d'histoire de la logique.
Les compléments comprennent des présentations alternatives et un
fragment de théorie de la démonstration. Les autres variantes sont
conservées dans les sources.
\par\medskip
Texte original : \href{https://openlogicproject.org/}{Open Logic Project}.
Traduction française et révision éditoriale produites par IA.
\par\medskip
Texte original et traduction : licence
\href{https://creativecommons.org/licenses/by/4.0/}{Creative Commons
Attribution 4.0 International}. Les notices de droits propres aux
composants sont conservées avec les sources.
\end{titlingpage}
\tableofcontents*
\chapter*{À propos du projet Open Logic}
\addcontentsline{toc}{chapter}{À propos du projet Open Logic}

Le \textit{Open Logic Text} est un manuel libre et collaboratif de
métalogique et de méthodes formelles. Il commence à un niveau
intermédiaire, après un cours d'introduction à la logique formelle.
Il s'adresse à des lecteurs sans formation spécialisée en
mathématiques, notamment aux étudiants en philosophie et en
informatique. L'exposé reste rigoureux.

La couverture et le niveau de détail varient selon les sujets.
Si vous trouvez une erreur ou avez une suggestion,
\href{https://github.com/OpenLogicProject/OpenLogic/wiki/Contributing}{signalez-la
à l'équipe du projet}.

Le projet fonctionne dans l'esprit des logiciels libres. Le texte
est accessible gratuitement et ses sources LaTeX sont diffusées sous
licence Creative Commons Attribution. Cette licence autorise chacun à
télécharger, utiliser, modifier, réorganiser, convertir et redistribuer
l'ouvrage, à condition d'en attribuer correctement la paternité.
Pour plus de renseignements, consultez le site du projet Open Logic :
\href{http://openlogicproject.org/}{openlogicproject.org}.

\mainmatter

\olimport[content/sets-functions-relations]{sets-functions-relations-complete}
\olimport[content/propositional-logic]{propositional-logic}
\olimport[content/first-order-logic]{first-order-logic}
\olimport[content/model-theory]{model-theory}
\olimport[content/computability]{computability}
\olimport[content/turing-machines]{turing-machines}
\olimport[content/incompleteness]{incompleteness}
\olimport[content/second-order-logic]{second-order-logic}
\olimport[content/lambda-calculus]{lambda-calculus}
\olimport[content/many-valued-logic]{many-valued-logic}
\olimport[content/normal-modal-logic]{normal-modal-logic}
\olimport[content/applied-modal-logic]{applied-modal-logic}
\olimport[content/intuitionistic-logic]{intuitionistic-logic}
\olimport[content/counterfactuals]{counterfactuals}
\olimport[content/set-theory]{set-theory}
\olimport[content/methods]{methods}
\olimport[content/history]{history}
\olimport[content/reference]{reference}
\olimport[content/proof-theory]{proof-theory}
\olimport[content/supplements]{supplements}

\backmatter
% Ces deux traductions sont citées dans les notices bibliographiques.
\nocite{Frege1953,Peter1967}
\bibliographystyle{plainnat}
\bibliography{bib/open-logic}
\end{document}
